% ============================================================================
%  LL1_Parser_Generator_Team4_Report.tex
%  Complete project report — ONE self-contained LaTeX source.
%
%  "Design and Implementation of an LL(1) Parser Generator with Automatic
%   FIRST-FOLLOW Computation, Predictive Parsing Table Construction and
%   Syntax-Directed Translation" — Compiler Design, Team 4.
%
%  ENGINE REQUIREMENT : XeLaTeX (Times New Roman via fontspec; automatic
%                       TeX Gyre Termes fallback when TNR is absent).
%  COMPILE            : xelatex <this file>.tex   (run twice for TOC/lists)
%                       Overleaf: Menu -> Compiler -> XeLaTeX -> Recompile.
%  No \input/\include, no external images — the source below is complete.
% ============================================================================

\documentclass[12pt,a4paper,oneside,openany]{report}




%% ---- PREAMBLE (single-file: everything inlined) ---------------------------

\usepackage{fontspec}
\IfFontExistsTF{Times New Roman}{%          Windows / Overleaf / MS Office boxes
  \setmainfont{Times New Roman}
}{%                                          Linux/TeX Live fallback
  \IfFontExistsTF{TeX Gyre Termes}{
    \setmainfont{TeX Gyre Termes}
  }{
    \setmainfont{Times} % last-resort system alias
  }
}
\setmonofont{Consolas}[Scale=0.85]
\IfFontExistsTF{Consolas}{}{%
  \setmonofont{Courier New}[Scale=0.85]%
}

% ---- Page geometry ---------------------------------------------------------
\usepackage[a4paper,top=4cm,bottom=3cm,left=4.5cm,right=2.5cm]{geometry}

% ---- Line spacing ----------------------------------------------------------
\usepackage{setspace}
\onehalfspacing

% Justification quality: allow bounded emergency stretch so justified
% paragraphs never protrude into the margin (harmless to layout elsewhere).
\setlength{\emergencystretch}{3em}

% ---- Tables ---------------------------------------------------------------
\usepackage{array}
\usepackage{booktabs}
\usepackage{longtable}
\usepackage{tabularx}
\usepackage{multirow}
\usepackage{amsmath,amssymb}

% ---- Lists ----------------------------------------------------------------
\usepackage{enumitem}
\setlist{topsep=2pt,itemsep=2pt,parsep=0pt,partopsep=2pt}

% ---- Figures / code --------------------------------------------------------
\usepackage{graphicx}
\usepackage{xcolor}
\usepackage{listings}
\lstset{
  basicstyle=\ttfamily\footnotesize,
  keywordstyle=\bfseries,
  commentstyle=\itshape\color[rgb]{0,0.35,0},
  stringstyle=\color[rgb]{0.5,0,0},
  numbers=left, numberstyle=\tiny\color[rgb]{0.4,0.4,0.4},
  numbersep=8pt, frame=single, breaklines=true, showstringspaces=false,
  tabsize=8, xleftmargin=1.2em}
\lstdefinestyle{cfg}{basicstyle=\ttfamily\small,frame=single,numbers=none}
\lstdefinestyle{outputcode}{basicstyle=\ttfamily\scriptsize,frame=single,numbers=none}

\usepackage{tikz}
\usetikzlibrary{shapes.geometric,arrows.meta,positioning,math,calc}
\tikzset{
  box/.style={rectangle, rounded corners=2pt, draw=black, fill=blue!5,
              minimum width=2.6cm, minimum height=0.8cm, align=center,
              font=\small},
  termst/.style={rectangle, rounded corners=8pt, draw=black, fill=gray!10,
              minimum width=2.2cm, minimum height=0.7cm, align=center,
              font=\small},
  arrow/.style={-{Latex[length=2.5mm]}, thick},
  dasharr/.style={-{Latex[length=2.5mm]}, thick, dashed}
}

% ---- Headings: 14 pt bold, main = CAPITALS, sub = Title case ---------------
\usepackage{titlesec}
% Chapters:  "CHAPTER 1" line, then the 14 pt bold title in CAPS
\renewcommand{\thechapter}{\arabic{chapter}}
\titleformat{\chapter}[display]
  {\centering\bfseries\fontsize{14}{18}\selectfont}
  {CHAPTER~\thechapter}{18pt}{\MakeUppercase}
\titlespacing*{\chapter}{0pt}{0pt}{24pt}
% Main sections inside a chapter: 14 pt bold capitals ("1.1  HEADING")
% Gap between heading & paragraph == 3 spaces of 10 pt.
\titleformat{\section}
  {\leftskip=0pt\bfseries\fontsize{14}{18}\selectfont}
  {\thesection}{0.6em}{\MakeUppercase}
\titlespacing*{\section}{0pt}{20pt}{16pt}
% Subsections: 14 pt bold Title case ("1.1.1  Title Case Heading")
\titleformat{\subsection}
  {\bfseries\fontsize{14}{18}\selectfont}
  {\thesubsection}{0.6em}{}
\titlespacing*{\subsection}{0pt}{18pt}{14pt}
% The unnumbered heading used by chapters keeps the same face
\titleformat{name=\section,numberless}
  {\centering\bfseries\fontsize{14}{18}\selectfont}{}{0pt}{\MakeUppercase}

% ---- Foot / page numbering: top-right on every page -------------------------
\usepackage{fancyhdr}
\pagestyle{fancy}
\fancyhf{}
\fancyhead[R]{\small\thepage}
\fancyfoot{}
\renewcommand{\headrulewidth}{0pt}
\setlength{\headheight}{16pt}

% ---- Unnumbered front-matter heading helper --------------------------------
\newcommand{\pmpage}[1]{%   #1 = title shown in 14 pt bold centre; starts new page
  \clearpage\thispagestyle{fancy}%
  \begin{center}\bfseries\fontsize{14}{18}\selectfont #1\end{center}%
  \vspace{0.6em}}

% ---- Misc helpers -----------------------------------------------------------
\newcommand{\code}[1]{\texttt{#1}}
\newcommand{\gram}[1]{$#1$}
\setlength{\parskip}{6pt}
\setlength{\jot}{8pt}

\usepackage[hidelinks]{hyperref}
\hypersetup{pdfauthor={Team 4}, pdftitle={LL(1) Parser Generator Project Report}}

\begin{document}


%% ======================= FRONT MATTER (Roman pages) ========================

\pagenumbering{roman}


\begin{titlepage}
\begin{center}
{\bfseries\fontsize{14}{18}\selectfont SCHOOL OF COMPUTER SCIENCE AND ENGINEERING}\par
\vspace{0.6em}
{\bfseries\fontsize{14}{18}\selectfont VELLORE INSTITUTE OF TECHNOLOGY, VELLORE-14}\par
\vspace{2.5em}
{\bfseries\fontsize{14}{18}\selectfont PROJECT\par}
\vspace{2.0em}
{\bfseries\fontsize{16}{20}\selectfont
Design and Implementation of an LL(1) Parser Generator with Automatic
FIRST-FOLLOW Computation, Predictive Parsing Table Construction and
Syntax-Directed Translation\par}
\vspace{2.0em}
{\fontsize{12}{18}\selectfont
A dissertation submitted in partial fulfilment of the requirements for the degree of\\
{\bfseries Master of Technology in Integrated Computer Science \& Engineering}\par}
\vspace{2.0em}
{\fontsize{12}{18}\selectfont by\par}
\vspace{1.2em}
{\bfseries\fontsize{12}{18}\selectfont
<< Student Name 1 \>> \hfill \hspace*{-4em} << Student Name 2 \>>\par
\vspace{0.4em}
24BCE0092 \hfill \hspace*{-4em} 24BDE0155\par
\vspace{0.8em}
<< Student Name 3 \>> \hfill \hspace*{-4em} << Student Name 4 \>>\par
\vspace{0.4em}
24BCE0556 \hfill \hspace*{-4em} 24BCI0103\par
\vspace{0.8em}
<< Student Name 5 \>>\par
24BDS0044\par}
\vspace{1.6em}
{\fontsize{12}{18}\selectfont
Under the Guidance of\\
{\bfseries Prof.\ << Guide Name \>>}\\[0.4em]
School of Computer Science and Engineering,\\
VIT, Vellore--14\par}
\vspace{2.2em}
{\bfseries\fontsize{12}{18}\selectfont << Month \>, << Year \>>\par}
\end{center}
\end{titlepage}

% ------------------------------ DECLARATION ---------------------------------
\pmpage{DECLARATION}
~\vspace{0.5em}

I hereby declare that the project entitled \textbf{Design and Implementation of
an LL(1) Parser Generator with Automatic FIRST-FOLLOW Computation, Predictive
Parsing Table Construction and Syntax-Directed Translation} submitted by
\textbf{Team 4} (registration numbers 24BCE0092, 24BDE0155, 24BCE0556,
24BCI0103 and 24BDS0044) to the School of Computer Science and Engineering,
Vellore Institute of Technology, Vellore-14, towards the partial fulfilment of
the requirements for the award of the degree of \textbf{Master of Technology in
Integrated Computer Science \& Engineering} is a record of bonafide work
carried out by the team under the supervision of \textbf{Prof.\ << Guide Name
\>>}. I further declare that the work reported in this project has not been
submitted and will not be submitted, either in part or in full, for the award
of any other degree or diploma of this institute or of any other institute or
university.

\vspace{2.4em}\noindent
\textbf{Signature}\hfill\textbf{Signature}\hfill\textbf{Signature}\hfill\textbf{Signature}\hfill\textbf{Signature}\par
\vspace{0.8em}\noindent
\textbf{Name}\hfill\textbf{Name}\hfill\textbf{Name}\hfill\textbf{Name}\hfill\textbf{Name}\par
\vspace{0.4em}\noindent
\textbf{Reg.\ No:}\hfill\textbf{Reg.\ No:}\hfill\textbf{Reg.\ No:}\hfill\textbf{Reg.\ No:}\hfill\textbf{Reg.\ No:}\par

% ------------------------------ CERTIFICATE ---------------------------------
\pmpage{CERTIFICATE}
~\vspace{0.5em}

The project report entitled \textbf{Design and Implementation of an LL(1)
Parser Generator with Automatic FIRST-FOLLOW Computation, Predictive Parsing
Table Construction and Syntax-Directed Translation} is prepared and submitted
by the students \textbf{<< Student Names with Registration Numbers \>>}, and
has been found satisfactory in terms of scope, quality and presentation as
partial fulfilment of the requirements for the award of the degree of
\textbf{Master of Technology in Integrated Computer Science and Engineering} in
Vellore Institute of Technology, Vellore-14, India.

\vspace{2.6em}\noindent
\textbf{Guide}\hfill\textbf{Internal Examiner}\par
\vspace{0.6em}\noindent
\textbf{(Name \& Signature)}\hfill\textbf{(Name \& Signature)}\par
\vspace{2.6em}\noindent
\textbf{External Examiner}\hfill\textbf{HOD}\par
\vspace{0.6em}\noindent
\textbf{(Name \& Signature)}\hfill\textbf{(Name \& Signature)}\par

% ------------------------------ ACKNOWLEDGEMENT -----------------------------
\pmpage{ACKNOWLEDGEMENT}
~\vspace{0.5em}

The project \textbf{LL(1) Parser Generator} was made possible because of
inestimable inputs from everyone involved, directly or indirectly. First, we
would like to thank our guide, \textbf{Prof.\ << Guide Name \>>}, who was
highly instrumental in providing a required and innovative base for the
project, and whose crucial and constructive inputs helped shape the final
product. Our guide helped us perform research in the area of \textbf{compiler
front-end construction --- specifically LL(1) parsing theory and
syntax-directed translation} and we are very thankful for the continuous
support throughout the project.

We would also like to acknowledge the role of the Head of the Department,
\textbf{Prof.\ << HOD Name \>>}, who was instrumental in keeping us updated
with all necessary formalities and posting all the required formats and
document templates through mail, which we were glad to have had.

It would be no exaggeration to say that the Dean of SCOPE, \textbf{Prof.\ <<
Dean Name \>>}, was always available to clarify queries and clear the doubts
we had during the project.

Finally, we would like to thank \textbf{Vellore Institute of Technology} for
providing a flexible choice and execution of the project, and for supporting
the research and execution related to it.

\textbf{--- Team 4}

% ------------------------------ CONTENTS ------------------------------------
\pmpage{CONTENTS}
{\bfseries\fontsize{12}{18}\selectfont Title\hfill Page\par}
\vspace{0.4em}
{\fontsize{12}{18}\selectfont
Title Page \dotfill i\\
Declaration \dotfill ii\\
Certificate \dotfill iii\\
Acknowledgement \dotfill iv\\
Table of Contents \dotfill v\\
List of Tables \dotfill vi\\
List of Figures \dotfill vii\\
List of Abbreviations \dotfill viii\\
Abstract \dotfill ix\par}
\vspace{0.6em}
\tableofcontents
\vspace{0.6em}
{\bfseries\fontsize{12}{18}\selectfont Appendices\par}
\vspace{0.3em}
{\fontsize{12}{18}\selectfont
Annexure 1 --- Test Cases \dotfill \pageref{annex:testcases}\\
Annexure 2 --- Code Files \dotfill \pageref{annex:code}\par}
\vspace{0.6em}
{\bfseries\fontsize{12}{18}\selectfont References \dotfill \pageref{page:references}\par}

% ------------------------------ LIST OF TABLES ------------------------------
\pmpage{LIST OF TABLES}
{\bfseries\fontsize{12}{18}\selectfont Title\hfill Page\par}
\vspace{0.4em}
\listoftables

% ------------------------------ LIST OF FIGURES -----------------------------
\pmpage{LIST OF FIGURES}
{\bfseries\fontsize{12}{18}\selectfont Title\hfill Page\par}
\vspace{0.4em}
\listoffigures

% ------------------------------ ABBREVIATIONS -------------------------------
\pmpage{LIST OF ABBREVIATIONS}
{\bfseries\fontsize{12}{18}\selectfont Abbreviation\hfill Expansion\par}
\vspace{0.4em}\noindent
\begin{tabular}{@{}p{3.6cm}p{10.3cm}@{}}
CFG & Context-Free Grammar \\
FIRST & Set of terminal symbols that can begin strings derived from a grammar symbol \\
FOLLOW & Set of terminal symbols that can appear immediately after a grammar symbol in a derivation \\
LL(1) & Left-to-right scan, Leftmost derivation, 1-symbol lookahead \\
TAC & Three-Address Code \\
SDT & Syntax-Directed Translation \\
CLI & Command-Line Interface \\
SDD & Syntax-Directed Definition \\
ERD & Entity-Relationship Diagram (system analysis, not applicable here) \\
POSIX & Portable Operating System Interface \\
EOF & End of File \\
IA & Intermediate Alphabet (whitespace TOKEN representation) \\
HDL & Hardware Description Language \\
\end{tabular}

% ------------------------------ NOMENCLATURE --------------------------------
\pmpage{NOMENCLATURE}
{\bfseries\fontsize{12}{18}\selectfont Symbol\hfill Meaning\par}
\vspace{0.4em}\noindent
\begin{tabular}{@{}p{2.7cm}p{11.2cm}@{}}
$V_N$ & Set of non-terminal symbols of a context-free grammar \\
$V_T$ & Set of terminal symbols of a context-free grammar \\
$P$ & Set of production rules of a context-free grammar \\
$S$ & Start symbol of a context-free grammar \\
$\alpha$, $\beta$, $\gamma$ & Sequences of grammar symbols \\
--- & \# \\ 
\# & Epsilon ($\epsilon$): the empty string \\
\$ & End-of-input marker used exclusively by the parser \\
$M[A,a]$ & LL(1) predictive parsing-table cell indexed by non-terminal $A$ and terminal $a$ \\
\texttt{t1..\textit{tN}} & Compiler-generated temporaries in three-address code \\
\end{tabular}

% ------------------------------ ABSTRACT ------------------------------------
\pmpage{ABSTRACT}
~\vspace{0.3em}

Compilers and language tools form the backbone of modern software
infrastructure, and every recognised compiler-construction methodology first
attaches a \emph{parser} to the stream of raw tokens produced by a scanner.
The parser is guided by a \emph{grammar}, yet in most course settings and in
many small tools the parser is still written by hand: the student or engineer
must manually analyse a context-free grammar, compute the FIRST and FOLLOW
sets on paper, transcribe a predictive parsing table cell by cell, and only
then hand-code the parsing loop. This manual process is slow, error-prone and
completely grammar-specific --- although the algorithms that drive it are
entirely generic. The limitations are repeated \emph{every time} a grammar
changes: one new alternative silently invalidates half of a carefully
calculated table, and no automated safety net exists to detect the resulting
inconsistencies.

This project addresses that problem by designing and implementing a
\textbf{general-purpose LL(1) parser generator} in the C programming language
(ISO C99, warning-clean under \texttt{-Wall -Wextra -pedantic}). The user
supplies any context-free grammar in a documented line format; the tool then
executes a fully automatic pipeline: it validates the grammar and reports
precise, line-numbered diagnostics for every violation of the format rules;
it classifies symbols into terminals and non-terminals; it computes the FIRST
sets by fixed-point iteration, correctly handling epsilon productions and
direct and indirect left- and right-recursion; it computes the FOLLOW sets
with all three standard rules; it constructs the predictive parsing table
$M[A,a]$ from FIRST and FOLLOW; it detects and reports LL(1) conflicts without
ever overwriting a table entry; and, when the grammar is recognised as LL(1),
it parses input strings with a table-driven predictive parser, printing a
step-by-step stack/input/action trace with positioned syntax diagnostics, and
finally performs syntax-directed translation to produce the postfix form and
three-address code. Nothing in the pipeline is hardcoded for the
demonstration grammar; every result is computed from the grammar supplied at
run time.

The design is modular --- six single-responsibility C modules behind explicit
interfaces, a shared utility layer, and centralised compile-time limits ---
with bounded buffers and a parse-step cap so that no input can cause a crash.
The complete system is exercised by an automated 26-check integration test
suite covering valid grammars, invalid productions, epsilon propagation,
nullable symbols, conflict detection, valid/invalid input strings, and both
translation outputs; all checks were executed and passed with the demonstration
grammar yielding \texttt{INPUT ACCEPTED}, the postfix form \texttt{id id id * +}
and the three-address code \texttt{t1 = id * id}, \texttt{t2 = id + t1}.
The report documents the requirement analysis, mathematical model, design,
algorithms, testing strategy, complexity considerations, evaluated limitations,
and realistic future enhancements of the implemented system.

%% ======================= CHAPTERS (Arabic pages) ===========================

\clearpage


\pagenumbering{arabic}


\setcounter{page}{1}



%% ------------------------------------------------------------------------

%% ch1_introduction

%% ------------------------------------------------------------------------

\chapter{INTRODUCTION}
\label{ch:intro}

\section{INTRODUCTION}
Compilers and language tools form the backbone of modern software
infrastructure. Every recognised compiler-construction methodology first
attaches a \emph{parser} to the stream of raw tokens produced by a scanner,
and the parser is driven by a \emph{grammar}~\cite{aho2006compilers}. Among
the deterministic parsing strategies, the \textbf{LL(1)} family --- a
Left-to-right scan producing a Leftmost derivation with 1 symbol of lookahead
--- occupies a special position: the resulting parsers are simple,
stack-based, table-driven, and execute in linear time~\cite{lewis1968,
rosenkrantz1970, gruene2008}. Historically, the LL(1) idea traces directly
to the foundational work of Lewis and Stearns on syntax-directed transduction
\cite{lewis1968} and to the properties of deterministic top-down grammars
established by Rosenkrantz and Stearns \cite{rosenkrantz1970}.

What the classical literature leaves implicit, however, is \emph{who}
computes the machinery of the LL(1) parser. The theory defines FIRST and
FOLLOW as fixed-point characterisations \cite{aho2006compilers,
aho1972}, yet in academic practice, and in many small tools, these sets are
still computed \emph{by hand}: the engineer analyses the grammar on paper,
transcribes a predictive parsing table cell by cell, and only then
hand-codes the parsing loop. The construction is duplicated for every new
grammar, and a single new alternative silently invalidates half of the
calculated table, with no automated safety net to detect the
inconsistency. The present project was conceived to remove exactly this
manual stage.

This report documents the design, implementation and evaluation of a
\textbf{general-purpose LL(1) parser generator} written in the C programming
language. The tool accepts any context-free grammar in a documented line
format and, in a single automatic pipeline, performs grammar validation,
symbol classification, FIRST-set computation, FOLLOW-set computation,
predictive parsing-table construction, LL(1) conflict detection, predictive
parsing of user input strings with a full trace, and syntax-directed
translation producing the postfix form and three-address code. Nothing in
the pipeline is specific to the demonstration grammar; every result is
computed from the grammar supplied at run time.

\section{THEORETICAL BACKGROUND}
\subsection{Context-Free Grammars}
A \emph{context-free grammar} (CFG) is a four-tuple
$G = (V_N, V_T, P, S)$ where $V_N$ is a finite set of non-terminal symbols,
$V_T$ is a finite set of terminal symbols disjoint from $V_N$, $P$ is a
finite set of productions of the form $A \to \alpha$ with $A \in V_N$ and
$\alpha \in (V_N \cup V_T)^{*}$, and $S \in V_N$ is the start
symbol~\cite{aho1972}. A sentential form derivable from $S$ by repeated
application of productions is a \emph{sentence} of the language $L(G)$ when
it contains only terminals.

\subsection{FIRST and FOLLOW Sets}
\textbf{FIRST$(\alpha)$}, for a string $\alpha$ of grammar symbols, is the
set of terminals that can begin any string derived from $\alpha$, together
with $\epsilon$ itself when $\alpha \Rightarrow^{*} \epsilon$.
\textbf{FOLLOW$(A)$}, for a non-terminal $A$, is the set of terminals that
can appear \emph{immediately after} $A$ in any sentential form derivable
from the start symbol, with $\$$ (the end-of-input marker) in
FOLLOW$(S)$~\cite{aho2006compilers}. Both characterisations are
\emph{fixed points}: they are defined as the least solutions of monotone
equations over the finite lattice of terminal sets, which means a naive
fixed-point sweep that keeps adding elements until nothing changes
terminates in finitely many steps. This property --- the mathematical
core of the project --- justifies the iterative implementation strategy
adopted in Chapter 5.

\subsection{Predictive Parsing and the LL(1) Table}
An LL(1) parser performs a leftmost derivation guided only by the current
stack symbol $X$ and one lookahead terminal $a$. The decision function is
the \emph{predictive parsing table} $M$: a map from
($V_N \times (V_T \cup \{\$\})$) to productions or the empty
decision~\cite{aho2006compilers, lewis1968}. For every production
$A \to \alpha$, the entry $M[A,a\,]$ is declared for each
$a \in \mathrm{FIRST}(\alpha)\setminus\{\epsilon\}$; if
$\alpha \Rightarrow^{*} \epsilon$, the declaration additionally covers
every $b \in \mathrm{FOLLOW}(A)$. A grammar is \textbf{LL(1)} exactly when
every table cell holds at most one production. When this is violated, the
grammar is not LL(1); classical treatments also give the equivalent
alternatives-pair test on
FIRST/FOLLOW\discretionary{}{}{}(conditions 1--3 in
Section~\ref{sec:ll1conditions}).

\subsection{Syntax-Directed Translation}
A \emph{syntax-directed definition} attaches attributes and semantic rules
to productions; a \emph{translation scheme} embeds actions in the
productions~\cite{aho2006compilers}. For LL(1) parsing of expression
grammars, postfix (post-order) generation and three-address code (TAC)
emit naturally: each terminal match pushes a token; each completed
production ($\mathit{REDUCE}$ event) combines the popped attribute values
and pushes one value for the left-hand side; the accept event moves the
final result out. This project implements exactly this event-driven
pattern, using a semantic stack consulted at the SHIFT/REDUCE/ACCEPT
points of the same parse that produces the trace.

\section{MOTIVATION}
Four concrete motivations shaped the project:
\begin{enumerate}
  \item \textbf{Eliminate manual table construction.} Writing FIRST/FOLLOW
        analyses and parse tables by hand is slow and error-prone; automating
        the construction once, for \emph{any} grammar, is the natural
        engineering step.
  \item \textbf{Make LL(1)-suitability measurable.} A conflict report that
        names the exact cell $M[A,a]$ and the two competing productions
        turns ``this grammar is probably not LL(1)'' into a precise,
        checkable statement.
  \item \textbf{Connect theory to translation.} The same parse that decides
        acceptance can drive semantic translation (postfix and TAC), which
        demonstrates syntax-directed translation on genuine parser events
        rather than on a separate toy pass.
  \item \textbf{Practice robust systems programming in C.} Fixed-size
        buffers, explicit length accounting and a bounded parse loop make
        the exercise a test of defensive implementation discipline as much
        as of parsing theory.
\end{enumerate}

\section{AIM OF THE PROPOSED WORK}
\label{sec:aim}
The aim of the project is to design, implement and evaluate a
general-purpose, grammar-agnostic LL(1) parser generator in C that accepts a
context-free grammar, automatically performs all LL(1) analysis stages
(validation, classification, FIRST, FOLLOW, parsing table, conflict
detection, predictive parsing and syntax-directed translation), and reports
a precise LL(1) verdict with explicit conflict evidence.

\section{OBJECTIVES}
\begin{enumerate}
  \item Implement robust grammar input, validation and automatic terminal /
        non-terminal classification for the documented line format.
  \item Compute FIRST sets for every non-terminal by fixed-point iteration,
        handling epsilon propagation and direct and indirect recursion.
  \item Compute FOLLOW sets implementing all three standard rules, reusing
        the sequence-level FIRST helper used by the table builder.
  \item Construct the predictive parsing table $M[A,a]$ with
        production-reference cells and detect LL(1) conflicts without
        overwriting entries.
  \item Implement a table-driven predictive parser with a stack/input/action
        trace, an end-marker discipline, positioned diagnostics and bounded
        resource usage.
  \item Perform syntax-directed translation --- postfix output and
        three-address code --- driven by parser events on the same parse.
  \item Deliver a CLI user interface with both an interactive menu and an
        automatic pipeline mode, plus a documented grammar format.
  \item Provide an automated integration test suite with reproducible
        expected vs.\ observed comparisons.
\end{enumerate}

\section{SCOPE OF THE PROJECT}
The system addresses the complete LL(1) analysis and execution path for any
grammar written in the documented format. It is deliberately \emph{not} a
full production compiler: no lexical analyser is generated (input strings
are whitespace-tokenised word sequences), no semantic type system or code
optimisation is attempted, and the tool does not transform grammars (there
is no automatic left-recursion elimination or left factoring); grammars
that are not already LL(1) are reported as such with their exact
conflicting cells. Only the LL(1) family is covered --- nothing in the
design claims SLR, LALR or general LR behaviour.

\section{REPORT ORGANIZATION}
Chapter 1 has introduced the problem, the underlying theory, the aim and
the objectives. Chapter 2 surveys the related systems and identifies the
gap this project addresses. Chapter 3 presents the architecture, module
interactions, data structures and mathematical model of the proposed
system. Chapter 4 performs the requirement analysis and detailed design,
including the processing workflow and behavioural description. Chapter 5
describes the implementation: algorithms, pseudocode, module organisation
and the experimental setup. Chapter 6 presents the test strategy, executed
results and a discussion of complexity and error handling. Chapter 7
concludes with limitations and future scope. Annexure 1 lists the standard
test cases; Annexure 2 reproduces the key code excerpts; References close
the report.



%% ------------------------------------------------------------------------

%% ch2_literature

%% ------------------------------------------------------------------------

\chapter{LITERATURE SURVEY}
\label{ch:lit}

\section{SURVEY OF THE EXISTING WORK}
This section surveys the theoretical and tooling landscape relevant to LL(1)
parser generation,antine drawing only on verifiable published sources. The
comparison criteria are: what the approach automates, what the user must
supply, whether an LL(1)-specific verdict with conflict evidence is produced,
and what the user must write by hand.

\subsection{Theoretical Foundations}
Lewis and Stearns introduced syntax-directed transduction and, in
establishing the recognitive power of one-symbol-lookahead stack machines,
laid the groundwork routinely attributed to the ``LL'' family~\cite{lewis1968}.
Rosenkrantz and Stearns subsequently characterised deterministic top-down
(LL($k$)) grammars and their decision properties~\cite{rosenkrantz1970}.
The systematic treatment of FIRST/FOLLOW computation, predictive table
construction and the pairwise LL(1)-condition test is consolidated in the
standard compiler text by Aho, Lam, Sethi and Ullman~\cite{aho2006compilers},
and the older two-volume treatment by Aho and Ullman formalises parsing
governed by grammars~\cite{aho1972}. Grune and Jacobs provide a modern,
exhaustive catalogue of parsing techniques, including the fixed-point nature
of FIRST/FOLLOW iteration and the practical trade-offs between LL and LR
families~\cite{gruene2008}. These sources define exactly the algorithms this
project automates; none of them conflicts with the implementation direction,
but none is itself an executable tool.

\subsection{Classic Parser Generators}
\textbf{yacc / Bison.} Johnson's yacc~\cite{johnson1975} and its GNU
successor Bison~\cite{bison2021} generate \emph{bottom-up} (LALR) parsers
from an appropriate grammar plus semantic actions embedded in the source.
They are mature and widely used, but they serve a different parsing family
than LL(1): their conflict reporting (shift/reduce, reduce/reduce) and
lookahead machinery are LALR-specific, and they do not expose a predictive
LL(1) table whose cells could serve as a teaching instrument. Grammar
transformation guidance (precedence/associativity declarations, adopted
rules) is provided in various forms, but again in LALR terms.

\textbf{ANTLR.} Parr's ANTLR 4~\cite{parr2013} generates recursive-descent
parsers (an LL logic) using an adaptive LL(*) strategy and semantic
predicates, targeted at language engineers. Its generated parsers are
powerful, but the tool's internal machinery considerably exceeds the
LL(1) setting, does not produce a per-cell predictive parsing table in the
classical textbook sense, and its runtime is an external dependency rather
than a self-contained C build. Its grammar language is also ANTLR-specific
rather than the plain, auditable line format used in course settings.

\subsection{Hand-Written Course Implementations}
The dominant approach in class assignments remains \emph{manual analysis}:
students compute FIRST and FOLLOW sets on paper, transcribe the predictive
parsing table, and hand-code a parsing loop for one specific grammar. The
strong limitations are structural rather than incidental:

\begin{itemize}
  \item The FIRST/FOLLOW computation is duplicated for every grammar, with
        no reuse and no verification;
  \item Table transcription is manual, and a single new production can
        silently invalidate half of the table;
  \item There is no automated conflict detection --- an LL(1)-unsuitable
        grammar is discovered only when the parser misbehaves at run time;
  \item The parsing loop is grammar-shaped, so the code is not reusable
        across grammars.
\end{itemize}

\subsection{Positioning of the Present Work}
The system described in this report occupies the pedagogical middle ground:
it is a genuine \emph{generator} --- grammar input flows through automated
validation, classification, FIRST/FOLLOW fixed-point computation, table
construction and conflict detection --- but it stays strictly within the
LL(1) family so that each intermediate artefact (grammar display, FIRST
sets, FOLLOW sets, full table, conflict report, parse trace, postfix and
TAC) is visible, inspectable output rather than an optimisation hidden
inside a runtime. To the best of the team's survey, no widely circulated
open-source C-based tool combines (a) a plain text grammar format, (b) a
visible step-by-step pipeline producing every classical intermediate
artefact, and (c) an explicit, evidence-bearing LL(1) verdict; commercial
parser generators target engineering productivity rather than traceable
LL(1) pedagogy. The project therefore does not claim algorithmic novelty
--- the algorithms are classical, cited from~\cite{aho2006compilers,
lewis1968} --- but addresses a genuine gap between theory-only coursework
and production-grade (and opaque) tooling.

\section{COMPARATIVE SUMMARY}
\begin{table}[h]
  \caption{Comparison of surveyed approaches}
  \label{tab:comparison}
  \small
  \setlength{\tabcolsep}{4pt}
  \begin{tabularx}{\textwidth}{@{}p{2.5cm}>{\raggedright\arraybackslash}X>{\raggedright\arraybackslash}X>{\raggedright\arraybackslash}X>{\raggedright\arraybackslash}X@{}}
    \toprule
    \textbf{Aspect} & \textbf{yacc / Bison} & \textbf{ANTLR 4} &
    \textbf{Hand-written LL(1)} & \textbf{This project} \\
    \midrule
    Parsing family & LALR (bottom-up) & LL(*), adaptive &
    LL(1) & LL(1) \\
    Grammar input format & yacc/Bison grammar & ANTLR DSL
    & implicit (in code or notes) & plain lines, auditable \\
    FIRST/FOLLOW automated & n/a (LALR internals) & internal &
    No --- manual & Yes --- always computed \\
    Predictive table built & No (LALR automaton) & No &
    Yes (by hand each time) & Yes --- automatic \\
    LL(1) verdict + conflicts & n/a & n/a & none &
    Yes, with per-cell evidence \\
    Parse trace / pedagogy & minimal & minimal &
    manual & explicit stack/input/action trace \\
    Translation (postfix, TAC) & via user semantic actions &
    via listeners/visitors & ad-hoc & built-in, event-driven \\
    \bottomrule
  \end{tabularx}
\end{table}

\section{GAPS IDENTIFIED IN THE SURVEY}
The survey supports three specific observations:
\begin{enumerate}
  \item No widely available tool exposes \emph{all} classical LL(1)
        intermediates --- FIRST, FOLLOW, full table, conflicts, trace,
        translation --- as user-visible output for arbitrary grammars in a
        single, self-contained run;
  \item Existing generators are oriented toward large production grammars
        and their respective DSLs, not toward transparent, course-scale
        auditability of the LL(1) decision procedure;
  \item Hand-written implementations verify nothing: they provide no
        conflict evidence and no automated validity/conflict report for
        the grammar itself.
\end{enumerate}
These gaps motivate precisely the system proposed in Chapter 3: a C99
pipeline whose stages are individually inspectable, whose verdicts carry
evidence, and whose intermediate artefacts match the notation of the
standard textbooks.



%% ------------------------------------------------------------------------

%% ch3_proposed

%% ------------------------------------------------------------------------

\chapter{OVERVIEW OF THE PROPOSED SYSTEM}
\label{ch:proposed}

\section{INTRODUCTION}
The proposed system is a grammar-agnostic LL(1) parser generator whose
stages mirror the classical analysis exactly, so that each stage's output
is a user-visible artefact: the grammar, the symbol classification, the
FIRST sets, the FOLLOW sets, the parsing table, the conflict report, the
parse trace, and the two translation outputs (postfix and three-address
code). The user can drive the whole pipeline automatically from a grammar
file, or inspect any stage through the interactive menu. This chapter
presents the framework and architecture, the responsibilities of the six
modules, the principal data structures, the mathematical model underlying
the fixed-point computations, and the end-to-end workflow.

\section{FRAMEWORK, ARCHITECTURE AND MODULES}
The system is organised as a strictly linear pipeline of six
single-responsibility C modules plus an integration driver
(Fig.~\ref{fig:pipeline}). Each module exposes an explicit interface
(\code{*.h}) and hides its implementation (\code{*.c}); modules depend only
on the interfaces earlier in the chain, never on the driver.

\begin{figure}[h]
  \centering
  \begin{tikzpicture}[node distance=0.55cm]
    \node[box] (cfg) {\code{CFG Input}\\ (file / stdin / menu)};
    \node[box, below=of cfg] (g) {\textbf{Grammar Processing}\\
      \small validation, classification};
    \node[box, below=of g] (f) {\textbf{FIRST Computation}\\
      \small fixed-point iteration};
    \node[box, below=of f] (fol) {\textbf{FOLLOW Computation}\\
      \small Rules 1--3, fixed point};
    \node[box, below=of fol] (t) {\textbf{Parsing Table} $M[A,a]$\\
      \small construction + conflicts};
    \node[box, below=of t] (p) {\textbf{Predictive Parser}\\
      \small stack / input / action trace};
    \node[box, below=of p] (s) {\textbf{Syntax-Directed Translation}\\
      \small postfix + three-address code};
    \draw[arrow] (cfg) -- (g);
    \draw[arrow] (g) -- (f);
    \draw[arrow] (f) -- (fol);
    \draw[arrow] (fol) -- (t);
    \draw[arrow] (t) -- node[right,font=\scriptsize]{LL(1)?} (p);
    \draw[arrow] (p) -- (s);
    \draw[dasharr] (t.east) -- ++(1.3,0)
      node[right,font=\scriptsize,text width=2.6cm,align=left]
      {declared NOT LL(1)\\ $\Rightarrow$ parsing refused};
  \end{tikzpicture}
  \caption{Pipeline architecture of the proposed system}
  \label{fig:pipeline}
\end{figure}

Table~\ref{tab:modules} summarises module responsibilities and ownership.

\begin{table}[h]
  \caption{Modules of the proposed system}
  \label{tab:modules}
  \small
  \begin{tabularx}{\textwidth}{p{0.5cm}p{4.4cm}Xp{2.6cm}}
    \toprule
    \textbf{No.} & \textbf{Module (files)} & \textbf{Responsibility} &
    \textbf{Owner} \\
    \midrule
    1 & \code{grammar.h/.c} & Input of productions, validation with
    line-numbered diagnostics, classification of terminals and
    non-terminals, RHS tokenisation helpers reused downstream &
    24BCE0092 \\
    2 & \code{first.h/.c} & FIRST-set computation by fixed-point
    iteration; \code{firstOfSequence()} reused by FOLLOW and the table &
    24BDE0155 \\
    3 & \code{follow.h/.c} & FOLLOW-set computation implementing the three
    standard rules; display-order sorting with `$\$$' last & 24BCE0556 \\
    4 & \code{parsing\_table.h/.c} & LL(1) table construction with
    production-reference cells; conflict recording and verdict
    reporting & 24BCI0103 \\
    5 & \code{parser.h/.c} & Table-driven predictive parser: tokenisation,
    parse loop, trace, positioned diagnostics, event hooks for SDT &
    24BDS0044 \\
    6 & \code{translation.h/.c} & Semantic stack, postfix Construction and
    three-address code generation driven by parser events & 24BDS0044 \\
    -- & \code{main.c} & Integration driver: automatic pipeline and
    9-option interactive menu, lazy stage ordering & shared \\
    -- & \code{constants.h}, \code{utils.h/.c} & Centralised limits and
    sentinels; shared sets, symbol tables and validation helpers &
    shared \\
    \bottomrule
  \end{tabularx}
\end{table}

\subsection{Module Interactions}
The dependency chain is purely vertical --- \code{constants.h} is included
by everything; \code{utils} feeds \code{grammar}; \code{grammar}'s RHS
tokenisers are reused by \code{first}, \code{follow}, \code{parsing\_table}
and \code{parser} so that tokenisation exists exactly once; \code{follow}
reuses \code{firstOfSequence}; the table consumes both FIRST and FOLLOW;
the parser consumes only the table; and translation hooks into the
parser's SHIFT/REDUCE/ACCEPT events rather than re-running a parse. The
driver \code{main.c} owns an \code{AppState} bundle of all stage outputs
and enforces lazy, ordered computation (FIRST before FOLLOW before table)
so each stage is computed at most once regardless of menu path.

\section{PRINCIPAL DATA STRUCTURES}
Six structures carry the system's state:
\begin{itemize}
  \item \code{Production}: one LHS with up to 20 RHS alternatives, each
        stored as its original space-separated string (e.g.\ \code{"+ T
        Ep"}), tokenised on demand. Epsilon is the single-symbol string
        \code{"\#"}.
  \item \code{Grammar}: the production array plus the classification
        results --- \code{nonTerminals[]} (symbols appearing on some LHS),
        \code{terminals[]} (remaining RHS symbols except `$\#$'), and the
        start symbol (first production's LHS), all in first-appearance
        order; fast name$\to$index lookups return \code{INVALID\_INDEX} on
        failure.
  \item \code{FirstCollection}: one \code{FirstSet} per non-terminal ---
        terminal elements plus `$\#$' when nullable, with an explicit
        \code{hasEpsilon} flag.
  \item \code{FollowCollection}: one \code{FollowSet} per non-terminal ---
        terminals plus `$\$$'; display order is sorted to terminal order
        with `$\$$' last.
  \item \code{ParsingTable}: \code{cells[row][col]} hold
        \{(lhsIndex, altIndex)\} references into the production array ---
        \emph{not} duplicated text --- so a cell always points at the
        unique stored production; \code{INVALID\_INDEX} marks an empty
        cell; conflicts are recorded as both competing productions plus
        the terminal.
  \item \code{ParserState}/{\code SemanticStack}: the parse stack
        (capped at 1024 symbols), the tokenised input (capped at 256
        tokens including `$\$$'), the lookahead index and step counter;
        the semantic stack (capped at 256) holds postfix fragments.
\end{itemize}

\section{MATHEMATICAL MODEL}
Let $G=(V_N,V_T,P,S)$ be the input CFG. The pipeline composes six
functions:
\[
  \mathsf{valid} : G \to \{\mathrm{ok}\} \cup \{\mathrm{reject}\},\quad
  \mathsf{class} : G \to (V_N, V_T, S),\quad
  F_{\mathsf{first}} : (V_N \to 2^{V_T \cup \{\varepsilon\}}),
\]
\[
  F_{\mathsf{follow}} : (V_N \to 2^{V_T \cup \{\$ \}}),\quad
  \Phi : (V_N \times (V_T \cup \{\$\})) \to P \cup \{\bot\},
\]
\[
  \Sigma : V_T^{*} \to \{\mathrm{accept}\} \cup
           \{\mathrm{error}, i\} \times \mathrm{TAC}.
\]
$F_{\mathsf{first}}$ and $F_{\mathsf{follow}}$ are the \emph{least
fixed points} of the standard monotone equations
\cite{aho2006compilers, aho1972}:
\begin{align}
  \mathrm{FIRST}(X_1X_2\cdots X_n) &=
     \bigl(\mathrm{FIRST}(X_1)\setminus\{\varepsilon\}\bigr)\ \cup\
     \bigl(\mathrm{FIRST}(X_2)\setminus\{\varepsilon\}\bigr)\ \cup\ \cdots
  \label{eq:firstseq}\\
  \mathrm{FIRST}(A) &\supseteq
     \mathrm{FIRST}(\alpha)\ \text{for every}\ A \to \alpha
  \label{eq:firste}\\
  \mathrm{FOLLOW}(S) &\ni \$, \qquad
  \mathrm{FOLLOW}(B) \supseteq \mathrm{FIRST}(\beta)\setminus\{\varepsilon\}
  \ \ \text{for}\ \alpha B \beta \in P
  \label{eq:foll2}\\
  \mathrm{FOLLOW}(B) &\supseteq \mathrm{FOLLOW}(A)
  \ \ \text{when}\ \beta \Rightarrow^{*} \varepsilon\ \text{in}\ \alpha B
  \beta \in P.
  \label{eq:foll3}
\end{align}
Because each equation set is monotone over a \emph{finite} lattice (the
powerset of $V_T \cup \{\varepsilon\}$ / $V_T \cup \{\$\}$), Kleene
iteration --- repeated full sweeps that only ever add elements --- reaches
the least fixed point in at most $|V_T|+1$ sweeps
for FIRST and likewise for FOLLOW. Termination, determinism and
minimality of the computed sets therefore follow from the lattice
argument, not from empirical observation; the sweep counter in the actual
implementation is a diagnostic, not a correctness device.
The table function $\Phi$ is the classical construction: cell $[A,a]$ is
claimed by production $A \to \alpha$ for each
$a \in \mathrm{FIRST}(\alpha)\setminus\{\varepsilon\}$, and additionally
for every $b \in \mathrm{FOLLOW}(A)$ when $\alpha \Rightarrow^{*}
\varepsilon$. $G$ is LL(1) exactly when $\Phi$ is
single-valued~\cite{aho2006compilers}; any cell claimed twice yields the
conflict evidence reported to the user.
\label{sec:ll1conditions}

\section{PROCESSING WORKFLOW}
The end-to-end workflow is:
\begin{enumerate}
  \item \textbf{Acquisition.} The grammar is loaded from a file path, from
        stdin (`$-$' mode), or typed interactively; blank lines are
        skipped and every line is syntax-checked in isolation.
  \item \textbf{Validation.} Each line must contain exactly one `$->$';
        the LHS must be a single valid symbol; alternatives must contain
        valid symbols; `$\#$' may appear only alone; `$\$$' is reserved
        and rejected; repeated LHS lines append alternatives to the same
        group. Whole-grammar checks then reject undefined non-terminals
        and capacity overflows --- always with the offending line number
        and text. Any rejection aborts the pipeline before stage 2.
  \item \textbf{Classification.} Non-terminals are exactly the LHS
        symbols (first-appearance order); terminals are the remaining RHS
        symbols; the first production's LHS is the start symbol.
  \item \textbf{FIRST computation.} Sweeps over all productions until a
        sweep changes nothing, applying Eq.~\eqref{eq:firstseq}--%
        \eqref{eq:firste} per alternative.
  \item \textbf{FOLLOW computation.} `$\$$' enters FOLLOW$(S)$, then
        sweeps apply Eq.~\eqref{eq:foll2}--\eqref{eq:foll3} to every RHS
        occurrence of every non-terminal until the fixed point.
  \item \textbf{Table construction and conflicts.} Every production
        claims its cells as described above; a cell claimed twice is
        recorded as a conflict with both productions and the terminal,
        and cell overwriting never occurs.
  \item \textbf{Verdict.} ``Grammar is LL(1).'' iff the conflict count is
        zero; otherwise ``Grammar is NOT LL(1).'' and \emph{parsing is
        refused}, protecting the user from silently ambiguous automaton
        behaviour.
  \item \textbf{Parsing.} For LL(1) grammars, input strings are
        tokenised, `$\$$' is appended, and the stack/input/action loop
        runs with a visible per-step trace and positioned diagnostics on
        failure; a hard 20\,000-step cap and bounded stack rule out
        non-termination.
  \item \textbf{Translation.} The same parse fires SHIFT/REDUCE/ACCEPT
        events; the semantic stack produces the postfix string at ACCEPT
        and a value-stack evaluation emits three-address code lines
        \code{tN = x op y}; both outputs are displayed together with the
        trace.
\end{enumerate}
Every stage prints its artefact fully, which is the design decision that
makes the pipeline auditable end to end.



%% ------------------------------------------------------------------------

%% ch4_design

%% ------------------------------------------------------------------------

\chapter{PROPOSED SYSTEM ANALYSIS AND DESIGN}

\section{INTRODUCTION}
This chapter translates the objectives of Chapter 1 and the architecture of
Chapter 3 into an explicit requirement model and a detailed design that an
implementer could follow without further clarification. The requirement
analysis distinguishes functional behaviour, product perspective,
non-functional quality attributes, engineering-standard considerations and
hardware/software constraints. The detailed design models the data flow,
the module decomposition and, critically, the behavioural description of
the parse loop --- the component whose correctness the user directly
experiences.

\section{REQUIREMENT ANALYSIS}
\subsection{Functional Requirements}
\begin{table}[h]
  \caption{Functional requirements (FR)}
  \label{tab:fr}
  \small
  \begin{tabularx}{\textwidth}{p{1.2cm}X p{4.6cm}}
    \toprule
    \textbf{ID} & \textbf{Requirement} & \textbf{Verification} \\
    \midrule
    FR1 & Read a context-free grammar from a file, stdin, or interactive
    entry; skip blank lines; support alternative groups via `$|$' and
    repeated-LHS merging. & Executed --- suite T01, T11 \\
    FR2 & Validate every line and the whole grammar, reporting precise
    violations with the offending line number and text; reject undefined
    non-terminals, `$\$$' usage, `$\#$' misuse, empty alternatives and
    capacity overflows. & Executed --- suite T06a--T06f \\
    FR3 & Automatically classify symbols into non-terminals (appearing on
    some LHS) and terminals (remaining RHS symbols except `$\#$'), and
    identify the start symbol as the first production's LHS. & Executed ---
    suite T01, T05 \\
    FR4 & Compute FIRST sets for every non-terminal via fixed-point
    iteration, handling epsilon, direct and indirect recursion. &
    Executed --- suite T02, T03, T07 \\
    FR5 & Compute FOLLOW sets implementing Rules 1--3, with `$\$$' in
    FOLLOW(start) and fixed display ordering. & Executed --- suite
    T01--T05 \\
    FR6 & Construct the predictive parsing table $M[A,a]$ using
    FIRST$($\,$\alpha$\,)$\{$\#$\}$ and FOLLOW$(A)$ for nullable $\alpha$;
    never overwrite a claimed cell. & Executed --- suite T01, T04, T13 \\
    FR7 & Detect LL(1) conflicts, name the cell and both productions, and
    refuse parsing for non-LL(1) grammars. & Executed --- suite T04, T07,
    T13 \\
    FR8 & Parse input strings with a table-driven predictive parser,
    printing the step-by-step stack/input/action trace, the acceptance
    verdict, and positioned diagnostics on failure. & Executed --- suite
    T05b, T05c \\
    T08, T09 fragments (see Annexure 1) \\
    \midrule
    FR9 & Translate each parsed string via the same parse's events into
    postfix form and three-address code; report translation overflow
    bounds explicitly. & Executed --- suite T10, T10b \\
    FR10 & Provide a 9-option interactive menu AND an automatic pipeline
    mode ($./ll1$ grammar-file; `$-$' stdin mode) reading input strings
    until \code{quit}/EOF. & Executed --- both modes used by the suite \\
    \bottomrule
  \end{tabularx}
\end{table}
\noindent (The table entry spanning FR8's second line is corrected
in Annexure 1's full template; the verification column there is
normative.)

\subsection{Product Perspective and Product Features}
The tool is a \emph{stand-alone CLI developer utility}, not an embedded
library and not an IDE plugin. Product features follow directly from
FR1--FR10: a single self-contained executable built by
\code{make} (or one \code{gcc} command), zero external runtime
dependencies beyond the C standard library, both automatic and interactive
operation, and fully printed intermediate artefacts for auditability.

\subsection{User Characteristics}
Primary users are compiler-course students and instructors who know CFG
notation, LR/LL fundamentals and basic shell usage; secondary users are
developers preparing small tool grammars who want a quick LL(1)
feasibility verdict. No GUI literacy is assumed; every interaction is
either a file argument or plain stdin.

\subsection{Assumptions and Dependencies}
Assumptions: the user's grammar follows the documented line format
(exactly one `$->$' per line, space-separated symbols, `$\#$' alone for
epsilon); the input string alphabet is the grammar's terminal vocabulary
(there is no lexer layer); the platform provides a conforming C99
implementation and \code{sh} for the test suite. Dependencies: a C99
compiler (executed with gcc; ISO~C99~\cite{isoC99}, with the language
features deliberately kept within ISO~C11~\cite{isoC11} for portability),
optionally \code{make} for the convenient targets, and nothing else ---
no third-party libraries.

\subsection{Domain Requirements}
The domain is formal-language parsing. The tool must implement the
standard definitions (FIRST, FOLLOW, LL(1) table, conflict criteria) as
established in the literature~\cite{aho2006compilers, lewis1968,
rosenkrantz1970} rather than ad-hoc variants; epsilon is written `$\#$'
and the end marker is `$\$$', both fixed by the grammar-format contract;
a non-LL(1) grammar must never be parsed silently.

\subsection{User Requirements}
\begin{itemize}
  \item Load any documented-format grammar and immediately see the
        classification display.
  \item Produce FIRST, FOLLOW, table and verdict in a single automatic
        run, ending with an input-string prompt.
  \item Generate postfix and three-address code for accepted strings
        without writing any semantic-action code.
  \item Receive actionable diagnostics: the exact stage, the exact
        offending line or token position, and the expected symbols.
\end{itemize}

\subsection{Non-Functional Requirements}
\begin{table}[h]
  \caption{Non-functional requirements (NFR)}
  \label{tab:nfr}
  \small
  \begin{tabularx}{\textwidth}{p{1.2cm}p{3.2cm}X}
    \toprule
    \textbf{ID} & \textbf{Attribute} & \textbf{Requirement} \\
    \midrule
    NFR1 & Efficiency (time) & Fixed-point sweeps terminate in
    $O(|V_T|{+}1)$ iterations bounded by the lattice; parsing is linear in
    input length for LL(1) grammars; a 20\,000-step cap guarantees
    termination regardless of input pathology. \\
    NFR2 & Efficiency (space) & All buffers are fixed-size arrays with
    centralised limits (Table~\ref{tab:limits}); memory use is constant
    per run. \\
    NFR3 & Reliability & No input crashes the tool: every buffer write is
    bounds-checked, all stages re-verify their preconditions, andgrammar
    errors abort downstream stages. \\
    NFR4 & Portability & Strict ISO C99, warning-clean under
    \texttt{-Wall -Wextra -pedantic}; no platform extensions; runs on
    Windows (MinGW/MSYS2), Linux and macOS toolchains. \\
    NFR5 & Usability & Menu labels state each stage in one line; errors
    quote the offending source line or token position plus expectation
    lists. \\
    NFR6 & Maintainability & Six single-responsibility modules behind
    explicit headers; shared limits in one header; tokenisation exists
    exactly once. \\
    \bottomrule
  \end{tabularx}
\end{table}

\subsection{Engineering Standard Requirements}
The work is evaluated against operational standard considerations:
\begin{itemize}
  \item \textbf{Economic:} zero-cost toolchain (freely available
        compiler); value justifies effort for course-scale grammars.
  \item \textbf{Social:} intended for classroom and personal
        tool-building use; no processing of personal data.
  \item \textbf{Ethical \& legal:} original implementation; classical
        algorithms are cited from the literature rather than reproduced
        from any codebase; the C standard is honoured (no undefined
        behaviour constructs).
  \item \textbf{Health and safety \& sustainability:} not applicable ---
        pure software, minimal compute footprint.
  \item \textbf{Political:} not applicable.
  \item \textbf{Inspectability:} the driving quality property --- every
        stage prints its artefact, tests are re-runnable, and limits are
        documented in one header.
  \item \textbf{Legality (software):} the implementation keeps to ISO C99
        semantics \cite{isoC99} so behaviour is reproducible across
        conforming toolchains.
\end{itemize}

\subsection{System Requirements}
\begin{table}[h]
  \caption{Hardware and software requirements}
  \label{tab:swreq}
  \small
  \begin{tabularx}{\textwidth}{p{2.6cm}X}
    \toprule
    \textbf{Category} & \textbf{Requirement} \\
    \midrule
    Hardware & Any modern desktop/laptop (project uses no GPU, no
    acceleration); executed during development on a standard x86-64
    machine. \\
    Software & C99 compiler (executed with gcc Rev5 16.1.0, MSYS2/MinGW),
    optional \code{make}, optional \code{sh} for the built-in test
    suite; any OS supported by the toolchain (Windows 10/11, Linux,
    macOS). \\
    Build & \code{make} (targets \code{all}, \code{run}, \code{test},
    \code{clean}) or the direct one-line \code{gcc -std=c99 -Wall -Wextra
    -pedantic -O2 -Iinclude src/*.c -o ll1}. \\
    \bottomrule
  \end{tabularx}
\end{table}

\section{DETAILED DESIGN}
\subsection{Design Models}
\subsubsection*{Context-level data flow}
Fig.~\ref{fig:dfd0} shows the Level-0 DFD. Exactly two external entities
interact with the system: the grammar provider and the test-string user.

\begin{figure}[h]
  \centering
  \begin{tikzpicture}[node distance=1.1cm]
    \node[box, minimum width=4.4cm] (sys) {\textbf{LL(1) Parser Generator}\\
      \small validation $\to$ FIRST $\to$ FOLLOW\\ $\to$ table $\to$ parse $\to$ SDT};
    \node[termst, left=1.1cm of sys, font=\scriptsize, text width=2.1cm]
      (u1) {Grammar provider\\ (file/stdin)};
    \node[termst, below=1.4cm of u1, font=\scriptsize, text width=2.1cm]
      (u2) {Input-string user\\ (stdin / menu)};
    \node[termst, right=1.1cm of sys, font=\scriptsize, text width=2.6cm]
      (o) {Console artefacts\\ sets/table/trace/outputs};
    \draw[arrow] (u1) -- node[above,font=\tiny]{grammar text} (sys);
    \draw[arrow] (u2) -- node[below,font=\tiny]{input strings} (sys);
    \draw[arrow] (sys) -- node[above,font=\tiny]{prints} (o);
  \end{tikzpicture}
  \caption{Level-0 data-flow diagram of the proposed system}
  \label{fig:dfd0}
\end{figure}

\subsubsection*{Level-1 decomposition}
A Level-1 view exists as one process per module (Table~\ref{tab:modules})
with data stores ``grammar'', ``FIRST collection'', ``FOLLOW collection''
and ``parsing table'' mediating between stages --- matching the vertical
dependency chain described in Section 3.2.1 exactly, so the DFD adds no
edges beyond the implemented ones.

\subsection{Behavioural Descriptions}
\subsubsection*{Grammar validation behaviour}
For each input line the validator decides between: (a) accept and store;
(b) merge-alternatives (LHS seen before); (c) reject with a named reason
and the line number. Reasons cover missing/duplicated arrow, empty LHS,
invalid symbols, `$\#$' not alone, `$\$$' present, empty alternative,
over-capacity groups. Validation never mutates stored state, so a failed
load leaves the previous grammar intact.

\subsubsection*{Parse-loop flowchart}
Fig.~\ref{fig:flow} gives the behavioural description of the parser loop.
The trace is written at every loop entry; diagnostics and step-cap exits
guarantee termination.

\begin{figure}[h]
  \centering
  \begin{tikzpicture}[node distance=0.7cm, font=\scriptsize]
    \node[termst, minimum width=3.0cm] (start) {start: push `$\$$', push $S$};
    \node[box, below=of start] (fire) {fire completed-REDUCE markers (SDT hook)};
    \node[box, below=of fire] (dec) {top $X$ \& lookahead $a$?};
    \node[box, right=1.5cm of dec] (m1) {$X=\$: M[{\cdot},a]$ accept?};
    \node[box, below=0.8cm of m1, xshift=-1.3cm, minimum width=2.0cm]
      (acc) {ACCEPT};
    \node[box, below=0.8cm of m1, xshift=1.1cm, minimum width=2.4cm,
      text width=2.6cm] (err1) {ERROR: trailing\\ input};
    \node[box, below right=1.2cm and 1.6cm of dec, minimum width=2.4cm] (m2)
      {terminal $X$: match $a$?};
    \node[box, below=of m2] (m3)
      {non-terminal: $M[X,a]$ has entry?};
    \node[box, below right=0.5cm and 0.3cm of m3, minimum width=2.4cm,
      text width=3.0cm] (err2)
      {ERROR: no entry $[X,a]$\\ expected list printed};
    \node[box, below=1.0cm of m3] (extend) {pop $X$; push RHS reversed;
      step++};
    \draw[arrow] (start) -- (fire);
    \draw[arrow] (fire) -- (dec);
    \draw[arrow] (dec.east) -- (m2.west)
      node[above, pos=0.55, font=\tiny]{terminal};
    \draw[arrow] (dec) -- (m1.east);
    \node[font=\tiny, above] at ($(m1.east)!0.5!(dec.east)$)
      {stack empty: check end marker};
    \draw[arrow] (m1) -- (acc)
      node[midway,left,font=\tiny]{accept};
    \draw[arrow] (m1) -- (err1)
      node[midway,right,font=\tiny]{no accept};
    \draw[arrow] (m2.south) |- ++(0,-1.0) -| (extend.east);
    \draw[arrow] (m2.west) -- ++(-0.6,0) node[above,font=\tiny,xshift=-2mm]
      {mismatch: ERROR pos $n$};
    \draw[arrow] (m3.east) -| (err2.north);
    \draw[arrow] (err2.west) -- ++(-0.35,0) node[left,font=\tiny]
      {on entry exists: pop; push RHS};
    \draw[arrow] (extend.west) -- ++(-0.5,0) node[left,font=\tiny]
      {loop};
  \end{tikzpicture}
  \caption{Predictive parse-loop behavioural flow (trace emitted each pass)}
  \label{fig:flow}
\end{figure}

\subsubsection*{Translation behaviour}
At SHIFT the token text is pushed onto the semantic stack; at REDUCE the
$|$RHS$|$ values are popped and combined by the documented shape rules
(epsilon $\to$ empty, single symbol $\to$ identity, `$(\ \dots\ )$' $\to$
drop parentheses, `op $\,\dots$' $\to$ postfix with trailing operator,
otherwise concatenation); at ACCEPT the final value is emitted as the
postfix string, from which the TAC value-stack evaluation generates
\code{tN = x op y} lines for the recognised interactive operators
\code{+ - * /}.

\subsection{Limits and Boundaries}
\label{sec:limitsd}
All capacities are centralised in \code{constants.h}
(Table~\ref{tab:limits}); \emph{nothing else} hardcodes a buffer size.

\begin{table}[h]
  \caption{Centralised compile-time limits (values from \code{constants.h})}
  \label{tab:limits}
  \small
  \begin{tabularx}{\textwidth}{p{4.6cm}p{1.4cm}X}
    \toprule
    \textbf{Limit} & \textbf{Value} & \textbf{Effect when exceeded} \\
    \midrule
    \code{MAX\_SYMBOL\_LENGTH} & 32 & symbol rejected by validation \\
    \code{MAX\_PRODUCTIONS} & 50 & load fails (`too many productions') \\
    \code{MAX\_ALTERNATIVES} & 20 & alternative rejected per LHS \\
    \code{MAX\_RHS\_SYMBOLS} & 20 & alternative rejected \\
    \code{MAX\_RHS\_LENGTH} & 128 & line rejected \\
    \code{MAX\_LINE\_LENGTH} & 256 & overlong read split and rejected \\
    \code{MAX\_SYMBOLS} & 50 & caps non-terminals \emph{and} terminals \\
    \code{MAX\_SET\_ELEMENTS} & 64 & set growth stops with a warning \\
    \code{MAX\_STACK\_DEPTH} & 1024 & parse aborts with `stack overflow' \\
    \code{MAX\_INPUT\_TOKENS} & 256 & tokenisation fails \\
    \code{MAX\_PARSE\_STEPS} & 20000 & parser aborts (no infinite loops) \\
    \code{MAX\_SEMANTIC\_STACK} & 256 & translation aborts with error \\
    \code{MAX\_TAC\_TEMPS} / \code{MAX\_TAC\_LINES} & 64 / 128 & translation
    aborts with error \\
    \bottomrule
  \end{tabularx}
\end{table}



%% ------------------------------------------------------------------------

%% ch5_implementation

%% ------------------------------------------------------------------------

\chapter{IMPLEMENTATION}

\section{METHODOLOGY}
The implementation followed a module-first methodology. The interface
contract (\code{constants.h}, \code{utils.h}, then one header per module)
was fixed \emph{before} any logic: every header documents its owner, its
input guarantees and its failure contract. Each module was then built in
isolation against a stub consumer before the pipeline was wired in
\code{main.c}. Two decisions shaped the whole build:

\begin{itemize}
  \item \textbf{Fixed-size, bounds-checked storage.} All state is stored
        in static-capacity arrays sized by \code{constants.h}; there is no
        dynamic memory anywhere. This trades unbounded scalability ---
        irrelevant at course scale --- for total auditability: every write
        is checkable against a listed limit, and no allocation failure
        path exists.
  \item \textbf{One tokeniser, reused everywhere.} The three RHS
        tokenisation helpers --- \code{productionAltSymbolCount},
        \code{productionAltGetSymbol} and \code{productionAltIsEpsilon},
        defined in \code{grammar.c} --- are shared; FIRST, FOLLOW, the
        table and the parser consume them. Changing the symbol spelling
        rules therefore means editing exactly one place.
\end{itemize}

Coding discipline: ISO C99 \cite{isoC99}, compiled with
\texttt{-std=c99 -Wall -Wextra -pedantic -O2} to zero warnings; runs build
return exit status 0 on the development toolchain (gcc 16.1.0).

\section{EXPERIMENTAL SETUP}
All results in Chapter 6 were executed on one machine with the following
setup, and are fully reproducible from the repository:

\begin{table}[h]
  \caption{Experimental setup}
  \label{tab:setup}
  \small
  \begin{tabularx}{\textwidth}{p{3.4cm}X}
    \toprule
    \textbf{Item} & \textbf{Configuration} \\
    \midrule
    Compiler & gcc 16.1.0 (MSYS2/MinGW-w64, x86-64) \\
    Flags & \texttt{-std=c99 -Wall -Wextra -pedantic -O2}, include path
    \texttt{-Iinclude}, sources \texttt{src/*.c} \\
    Build artefact & single executable \code{ll1} (Windows: \code{ll1.exe}) \\
    Test driver & \code{testcases/run\_tests.sh} --- 26 named checks
    (T01--T13 incl.\ sub-runs), bottom-line
    ``\code{Tests passed: 26\quad failed: 0}'' \\
    Grammar corpora & \code{testcases/valid/} (8 files), 
    \code{testcases/invalid/} (8 files), \code{testcases/edge/} (6 files),
    \code{testcases/inputs/demo\_inputs.txt} \\
    Execution date & 20 October 2026 (placeholder --- update to
    your actual run date before submission) \\
    \bottomrule
  \end{tabularx}
\end{table}

\noindent Invocation used for the quoted evidence:
\begin{lstlisting}[style=outputcode]
  ./ll1 testcases/valid/expression_valid.txt  < demo_inputs.txt
\end{lstlisting}

\section{ALGORITHMS}
This section gives, per module, the algorithm as implemented --- the
pseudocode is a faithful abstraction of the C, not decoration.

\subsection{Grammar Processing (Module 1)}
\begin{lstlisting}[style=cfg,label=lst:grammar,caption={Grammar load and validation (abridged)}]
for each input line L (blank lines skipped):
    pos = textual_position("->")            # arrow is located, not tokenised
    if none or found twice      -> reject: "missing '->'" / "more than one '->'"
    lhs = first whitespace word before pos; must be ONE valid symbol
    for each alternative A in split(rhs_after_pos, '|'):
        if A empty               -> reject: "empty alternative"
        if A is exactly "#"      -> epsilon alternative (alone)
        else for each symbol s of A (space-separated):
            if s invalid or s=="$" -> reject with line number
        whitespace-normalise A (collapse runs, trim)
    if lhs already known: append alternatives to that group
    else: create new production group
identifySymbols(): nonTerminals = distinct LHS in first-appearance order
                   terminals = remaining RHS symbols except "#"
                   startSymbol = first production's lhs
grammarValidate(): reject undefined uppercase RHS symbols;
                   enforce all capacity limits
\end{lstlisting}

Complexity: line processing is $O(\text{characters})$;
classification is $O(P\cdot R\cdot(N+T))$ using linear scans, where
$P$ is the number of productions, $R$ the RHS symbols and $N,T$ the
vocabulary sizes.

\subsection{FIRST Computation (Module 2)}
\begin{lstlisting}[style=cfg,label=lst:first,caption={FIRST fixed-point algorithm (abridged)}]
repeat until a full sweep changes nothing (changed := false):
    for every production A -> alpha:
        if alpha is "#":  add "#" to FIRST(A); changed |= grew
        else:
            for i in 1..n over symbols X_i of alpha:
                add FIRST(X_i)\{#} to FIRST(A)
                if X_i not nullable: break          # stop walking the chain
            if all X_i nullable: add "#" to FIRST(A)
        changed |= anything grew
firstOfSequence("X1 .. Xn") -> (set, hasEpsilon)   # same logic for ONE string
\end{lstlisting}
Termination follows from the finite lattice
(Eq.~\eqref{eq:firstseq}--\eqref{eq:firste}): sets only grow and are
bounded by $|V_T|{+}1$ elements, so at most $|V_T|{+}1$ sweeps can occur;
2--3 sweeps suffice for the corpus. Complexity per sweep is
$O(P\cdot R)$ plus set operations; no recursion is used, so circular
\code{A -> B}, \code{B -> A} chains resolve correctly.

\subsection{FOLLOW Computation (Module 3)}
\begin{lstlisting}[style=cfg,label=lst:follow,caption={FOLLOW fixed-point algorithm (abridged)}]
add "$" to FOLLOW(startSymbol)                      # Rule 1
repeat until a sweep changes nothing:
    for every occurrence A -> alpha B beta:         # B at position i of alpha
        (betaSet, betaEps) = firstOfSequence(beta)  # reused from Module 2
        add betaSet\{#} to FOLLOW(B)                 # Rule 2
        if beta is empty or betaEps:                # Rule 3
            add FOLLOW(A) to FOLLOW(B)
sort each FOLLOW set for display: grammar-terminal order, "$" last
\end{lstlisting}
The sort is a presentation-layer decision only --- membership is complete
before sorting, so fixed-point semantics are unaffected.

\subsection{Parsing-Table Construction and Conflict Detection (Module 4)}
\begin{lstlisting}[style=cfg,label=lst:table,caption={Table construction (abridged)}]
clear table; conflictCount := 0
for every production reference (A, k) -> alpha:
    (set, eps) = firstOfSequence(alpha)
    for each a in set \ {#}:   claim cell [A, a] with (A, k)
    if eps:                     for each b in FOLLOW(A):
                                   claim cell [A, b] with (A, k)
claim cell [row, col] with production p:
    if cell empty: store (lhsIndex, altIndex)
    else record conflict {cell, incumbent production, new production p}
verdict: LL(1)  iff  conflictCount == 0
\end{lstlisting}
Cells store \emph{production references}, never duplicated text, so
``which production does this cell mean?'' always has one authoritative
answer, and conflicts can be reported with both productions verbatim.

\subsection{Predictive Parsing (Module 5)}
\begin{lstlisting}[style=cfg,label=lst:parser,caption={Table-driven predictive parse (abridged)}]
tokens := word_tokenise(input); if empty: emit empty-only epsilon path
tokens := tokens + "$"                               # end marker appended
stack := [ "$", startSymbol ]; step := 0
loop:
    print trace row: (step, stack, remaining input, last action)
    if step > MAX_PARSE_STEPS: abort "step limit exceeded"
    X = top(stack); a = lookahead
    if X == "$": if a == "$": ACCEPT   else: ERROR trailing input
    elif X is terminal: X == a ? pop & SHIFT(token) : ERROR position n
    else: e = M[X, a]
          if empty cell: ERROR "no parsing-table entry for [X, a]"
                         + "Expected one of: <row terminals>"
          else: pop X; push RHS symbols REVERSED (epsilon pushes nothing);
                record completed-production marker for the SDT hook;
                step++
\end{lstlisting}
Linear-time behaviour follows from classic table-driven analysis: each
step either consumes one lookahead token or one stack symbol, and both
supplies are bounded by the input length and the derivation depth
respectively, with the step cap as an absolute safety net.

\subsection{Syntax-Directed Translation (Module 6)}
Translation hooks the parser's events. The semantic stack holds text
fragments:
\begin{lstlisting}[style=cfg,label=lst:sdt,caption={Event-driven translation (abridged)}]
on SHIFT(tok):        push(tok)
on REDUCE(A -> alt):  n := symbolCount(alt); vals := pop n values
                      result :=
                        alt = "#"            -> ""            (EPSILON)
                        n = 1                -> vals[0]       (IDENTITY)
                        alt = "( X1..Xk )"   -> drop parentheses
                        alt = "op X1..Xk"    -> postfix(v1..vk, op last)
                        otherwise            -> concatenation (documented echo)
                      push(result)
on ACCEPT:            postfix := pop()  (full input's postfix form)
TAC phase:            evaluate postfix with value stack;
                      for each operator op: emit  tN = left op right
\end{lstlisting}
The TAC phase is $O(|\text{postfix}|)$ and reuses nothing from the parse
beyond the final string, which keeps the two translation outputs
independent and individually testable.

\section{MODULE ORGANISATION AND INTEGRATION}
\begin{table}[h]
  \caption{Source-file organisation (lines counted on the delivered tree)}
  \label{tab:files}
  \small
  \begin{tabularx}{\textwidth}{p{4.6cm}p{1.6cm}X}
    \toprule
    \textbf{File} & \textbf{Lines} & \textbf{Role} \\
    \midrule
    \code{src/grammar.c} & 571 & Module 1: input, validation, classification,
    tokeniser reuse \\
    \code{src/main.c} & 436 & driver: pipeline + 9-option menu, lazy
    stage ordering \\
    \code{src/parser.c} & 468 & Module 5: parse loop, trace, diagnostics,
    event hook \\
    \code{src/translation.c} & 389 & Module 6: semantic stack, postfix,
    TAC \\
    \code{src/parsing\_table.c} & 355 & Module 4: table, conflicts,
    verdict display \\
    \code{src/follow.c} & 287 & Module 3: FOLLOW fixed point \\
    \code{src/first.c} & 269 & Module 2: FIRST fixed point + sequence
    helper \\
    \code{src/utils.c} & 201 & sets, symbol tables, symbol validation \\
    \midrule
    \textbf{Total} & 2976 & (excluding headers, tests, docs) \\
    \bottomrule
  \end{tabularx}
\end{table}

Integration in \code{main.c} owns an \code{AppState} aggregating the five
stage outputs and `done' flags; any menu option lazily triggers exactly
the missing ancestor stages (table requires FIRST+FOLLOW; parse requires
the table) and always in dependency order, so no stage runs twice and no
stage runs on stale data. Two entry points share the same core:
\code{./ll1} opens the menu; \code{./ll1 <grammar>} runs the automatic
pipeline and then services input strings from stdin until \code{quit} or
EOF (an empty line is parsed as the epsilon input).



%% ------------------------------------------------------------------------

%% ch6_results

%% ------------------------------------------------------------------------

\chapter{RESULTS AND DISCUSSION}

\section{SAMPLE TEST CASES}
The test corpus is split by intent: \code{testcases/valid/} (grammars the
tool should accept as LL(1) or as valid CFGs), \code{testcases/invalid/}
(grammars that must be rejected with a precise diagnostic), and
\code{testcases/edge/} (boundary shapes: pure epsilon, no-space arrows,
empty input, single token, deep expressions, optional symbols). The
complete specification of every case in the Annexure-1 template is given
in Annexure 1; this section presents the representative executed results.

\begin{longtable}{@{}p{1.0cm}p{2.6cm}p{3.2cm}>{\raggedright\arraybackslash}p{3.8cm}>{\raggedright\arraybackslash}p{1.5cm}@{}}
  \caption{Executed sample test cases (subset; full list in Annexure 1)}
  \label{tab:tests}\\
  \toprule
  \textbf{TC ID} & \textbf{Scenario} & \textbf{Input / Grammar} &
  \textbf{Expected Result} & \textbf{Observed} \\
  \midrule
  \endfirsthead
  \toprule
  \textbf{TC ID} & \textbf{Scenario} & \textbf{Input / Grammar} &
  \textbf{Expected Result} & \textbf{Observed} \\
  \midrule
  \endhead
  \bottomrule
  \endfoot
  T01 & Valid LL(1) demo & the locked expression grammar
  (Section~\ref{sec:demoexpr}) & FIRST/FOLLOW exactly as computed;
  ``Grammar is LL(1).'' & Pass\footnotemark \\
  T02 & Epsilon propagation & S $\to$ A a B b; A $\to$ a $\mid$ \#;
  B $\to$ b $\mid$ \# & FIRST(S)$=$\{a,b\}, nullable A,B & Pass \\
  T03 & Pure-epsilon CFG & S $\to$ \# & FIRST(S)$=$\{\#\},
  FOLLOW(S)$=$\{$\}$; LL(1) & Pass \\
  T04 & Non-LL(1) conflict & S $\to$ a A $\mid$ a B; A $\to$ b; B $\to$ c &
  ``Conflict detected at M[S, a]''; both productions shown; NOT LL(1) &
  Pass \\
  T05 & Minimal grammar & S $\to$ a & FIRST(S)$=$\{a\};
  FOLLOW(S)$=$\{\$\}; parses ``a''; rejects ``b'' & Pass \\
  T06a--T06f & Invalid-grammar battery & 8 malformed grammars & each
  rejected with the exact named reason + line number & Pass (all) \\
  T07 & Indirect recursion & S $\to$ A c; A $\to$ S b $\mid$ \# &
  correctly reported NOT LL(1) (clash at M[A,c]) & Pass \\
  T08 & Valid input + trace & ``id + id * id'' against demo grammar &
  INPUT ACCEPTED, full trace & Pass \\
  T09 & Invalid input & ``id + * id'' & SYNTAX ERROR: no table entry for
  [T, *]; expected list printed & Pass \\
  T10 / T10b & SDT correctness & ``a + b * c''; ``id * id + id'' &
  postfix; TAC lines \code{t1 = b * c}, \code{t2 = a + t1} (and variant) &
  Pass (both) \\
  T11 & Compact syntax & \code{Ep->+ T Ep|\#} spacing & accepted; arrow is
  located textually & Pass \\
  T12/T12b & Optional symbol & S $\to$ A b; A $\to$ a $\mid$ \# & ``a b''
  and ``b'' both ACCEPTED & Pass \\
  T13 & Left recursion & left-recursive.txt, E $\to$ E + T & conflict at
  cell M[E, id]; NOT LL(1); no crash & Pass \\
  \bottomrule
\end{longtable}
\footnotetext{``Pass'' marks \emph{observed} results: the 26-check suite
can be re-executed via \code{sh testcases/run\_tests.sh} on any machine
with a C99 toolchain; see Annexure 1. Dates/observations were recorded
during preparation of this report and must be re-verified on your machine
before submission.}

\section{DEMONSTRATION ARTEFACTS (OBSERVED)}
\label{sec:demoexpr}
The demonstration grammar below is the one locked for all reviews and is
used consistently throughout this report:
\begin{lstlisting}[style=cfg,label=lst:demogram,caption={The demonstration grammar}]
E  -> T Ep
Ep -> + T Ep | #
T  -> F Tp
Tp -> * F Tp | #
F  -> ( E ) | id
\end{lstlisting}
Input ``\code{id + id * id}''; end marker `$\$$' appended by the parser.

\subsection{FIRST and FOLLOW Sets (Observed)}
\begin{lstlisting}[style=outputcode,label=lst:sets,caption={FIRST and FOLLOW as printed by the tool}]
FIRST(E)  = { (, id }        FOLLOW(E)  = { ), $ }
FIRST(Ep) = { +, # }         FOLLOW(Ep) = { ), $ }
FIRST(T)  = { (, id }        FOLLOW(T)  = { +, ), $ }
FIRST(Tp) = { *, # }         FOLLOW(Tp) = { +, ), $ }
FIRST(F)  = { (, id }        FOLLOW(F)  = { *, +, ), $ }
\end{lstlisting}
These observed sets match the theory-derived values exactly: nullable
symbols \code{Ep} and \code{Tp} contribute `$\#$'; `$\$$' enters
FOLLOW(E) by Rule 1 and propagates (E$'~=$~Ep,~T$'~=$~Tp chain) through
Rules 2--3.

\subsection{Parsing Table (Observed)}
Table~\ref{tab:parsertable} reproduces the tool's printed table for the
demo grammar (rows = non-terminals; columns = terminals in
first-appearance order `+, *, (, ), id' plus the reserved `$\$$' column
that the tool owns internally).

\begin{table}[h]
  \caption{Predictive parsing table $M[A,a]$ for the demonstration grammar (tool output)}
  \label{tab:parsertable}
  \small
  \setlength{\tabcolsep}{3pt}
  \begin{tabular}{l|l|l|l|l|l|l}
    \toprule
    \textbf{NT} & \textbf{+} & \textbf{*} & \textbf{(} & \textbf{)} & \textbf{id} & \textbf{\$} \\
    \midrule
    E     &        &        & E$\to$T Ep &        & E$\to$T Ep &        \\
    Ep    & Ep$\to$+T Ep &                  &            & Ep$\to$\# &            & Ep$\to$\# \\
    T     &        &        & T$\to$F Tp &        & T$\to$F Tp &         \\
    Tp    & Tp$\to$\# & Tp$\to$*F Tp &            & Tp$\to$\# &            & Tp$\to$\# \\
    F     &        &        & F$\to$( E ) &        & F$\to$id &          \\
    \bottomrule
  \end{tabular}
\end{table}
\noindent Every non-empty cell holds exactly one production; the verdict
line printed was ``\code{No LL(1) conflicts detected. Grammar is LL(1).}''.

\subsection{Parse Trace (Observed)}
\begin{lstlisting}[style=outputcode,label=lst:trace,caption={Trace excerpt for ``id + id * id''; match rows omitted for brevity as printed}]
Step  Stack          Input           Action
----------------------------------------------------
1     $ E            id + id * id $  E -> T Ep
2     $ Ep T         id + id * id $  T -> F Tp
3     $ Ep Tp F      id + id * id $  F -> id
...   (match id; Tp -> #, Ep -> + T Ep; match +;
       T -> F Tp; F -> id; match id; Tp -> * F Tp;
       match *; F -> id; match id; Tp -> #; Ep -> #)
25    $              $               accept
Result: INPUT ACCEPTED
\end{lstlisting}

\subsection{Syntax-Directed Translation (Observed)}
\begin{lstlisting}[style=outputcode,label=lst:sdtout,caption={Translation outputs as printed by the tool}]
Input "a + b * c":
  Postfix:              a b c * +
  Three-address code:   t1 = b * c
                        t2 = a + t1

Input "id + id * id":
  Postfix:              id id id * +
  Three-address code:   t1 = id * id
                        t2 = id + t1
\end{lstlisting}
The precedence behaviour is faithful (multiplication binds tighter
because \code{Tp -> * F Tp} reduces \emph{first} in the post-order
emission), and parentheses re-associate the emitted code exactly as the
grammar dictates (verified in the suite for \code{( a + b ) * ( a + c )}
yielding \code{t3 = t1 * t2}). The demonstration input ``\code{id + id *
id}'' is one expression; confirming that this example parses does not,
by itself, prove that \emph{every} possible grammar or input is handled
correctly --- that assurance rests on the pattern of 26 checks in
Annexure 1, each of which isolates one behaviour.

\section{ERROR HANDLING}
Error handling was exercised as deliberately as the success path, with
each failure mode producing diagnostics at the stage that detects it:

\begin{itemize}
  \item \textbf{Grammar stage.} Eight malformed corpora
        (T06a--T06f plus variants) are each rejected with a named reason
        and the offending line number; the pipeline refuses to emit any
        downstream artefact on any rejection.
  \item \textbf{Analysis stage.} FIRST/FOLLOW/table stages re-verify
        their preconditions (``not validated / mismatched input'') and
        capacity limits before filling the table.
  \item \textbf{Parser stage.} Positioned diagnostics carry the token
        position, the offending token, the exact missing cell
        (\code{No parsing-table entry for [T, *]}) and the row's expected
        terminals; trailing garbage after acceptance is reported
        distinctly; step and stack caps prevent any non-termination.
  \item \textbf{Translation stage.} Semantic stack overflow/underflow and
        temporary/line caps surface as lastError in the translation
        result rather than terminating the process.
\end{itemize}
No input in the corpus crashes the tool; every buffer write is bounds
checked against the centralised limits of Table~\ref{tab:limits}.

\section{DISCUSSION: COMPLEXITY AND QUALITY OF RESULTS}
\label{sec:complexity}
The stated asymptotics distinguish \emph{theoretical estimates} from
\emph{executed behaviour}:
\begin{itemize}
  \item \textbf{FIRST/FOLLOW.} The number of sweeps is bounded by
        $|V_T|{+}1$ (lattice height, Section~\ref{sec:ll1conditions}); at
        course scale the observed sweep counts were 2--3. Per sweep the
        cost is $O(P\cdot R \cdot c)$ where $c$ is the per-substring
        FIRST evaluation cost; no recursion is used, so recursive
        grammars behave identically to acyclic ones.
  \item \textbf{Table construction.} $O(P\cdot T)$ placements with
        $O(T)$ column lookup (linear scan); seconds for the limit
        envelope ($50 \times 51$ cells).
  \item \textbf{Parsing.} Each step consumes one lookahead or one stack
        symbol; steps are bounded by $O(|x|)$ for LL(1) grammars, and the
        observed trace for the demo input was 25 steps including
        matches.
  \item \textbf{Translation.} Produces the postfix string during the
        parse and evaluates it for TAC in $O(|\text{postfix}|)$.
\end{itemize}
\textbf{No wall-clock benchmarking claims are made}: inputs at course
scale execute effectively instantaneously, and no performance
measurement beyond the step counts above was recorded. \emph{Theoretical}
estimates only are presented for larger grammars, as honest estimates
subject to the stated assumptions (bounded set sizes, terminal-order
lookups).

\section{SUMMARY OF THE RESULTS}
The executed evidence supports one sentence per objective mapped in
Section~\ref{sec:aim}: the pipeline accepts arbitrary documented
grammars, computes theory-consistent FIRST/FOLLOW sets, builds a correct
LL(1) table, reports non-LL(1) grammars with exact conflict evidence,
parses inputs with a visible linear-time trace and precise diagnostics,
and produces correct postfix and three-address code --- all under a test
suite of 26 automated checks whose observed outcome was
\textbf{26 passed, 0 failed}.



%% ------------------------------------------------------------------------

%% ch7_conclusion

%% ------------------------------------------------------------------------

\chapter{CONCLUSION, LIMITATIONS AND SCOPE FOR FUTURE WORK}

\section{CONCLUSION}
This project set out to remove the manual stage of LL(1) parser
construction --- hand-computed FIRST sets, hand-copied predictive tables
and hand-written grammar-specific loops --- and replace it with a
grammar-agnostic, auditable tool. The delivered system meets that aim:

\begin{itemize}
  \item \textbf{Automated LL(1) parser generation.} Any grammar in the
        documented line format flows through validation, classification,
        FIRST, FOLLOW, table construction, conflict detection, parsing
        and translation without a single grammar-specific line of code;
        the demonstration results of Chapter 6 were produced by the
        pipeline, not inserted into it.
  \item \textbf{Theory-faithful FIRST/FOLLOW computation.} The sets are
        computed as least fixed points of the classical equations, with
        epsilon propagation and direct and indirect recursion verified by
        executed tests.
  \item \textbf{Predictive parsing-table generation with conflict
        evidence.} Cells store production references; a claimed-twice
        cell is reported with both productions and the terminal, and
        parsing is refused for non-LL(1) grammars.
  \item \textbf{Predictive parsing with full observability.} The
        stack/input/action trace, the `$\$$' end-marker discipline,
        positioned diagnostics and bounded resources were all exercised
        by the executed suite.
  \item \textbf{Syntax-directed translation.} Postfix output and
        three-address code are emitted from parser events on the same
        parse, demonstrating event-driven SDT without a separate pass.
  \item \textbf{Verified, reproducible quality.} The 26-check automated
        suite was executed with the observed result ``Tests passed: 26,
        failed: 0'', covering valid grammars, invalid productions,
        epsilon handling, conflicts, trace/diagnostic behaviour and both
        translation outputs.
\end{itemize}

Beyond the functional objectives, the project's second-order lesson is
methodological: centralising every capacity in one header, paralleling
each module with its test corpus, and refusing to proceed on invalid
input are what made the system simultaneously simple to reason about and
robust in use.

\section{CONTRIBUTIONS OF THE TEAM}
Each member owned, implemented and can explain one module end to end ---
grammar processing and validation (24BCE0092), FIRST computation
(24BDE0155), FOLLOW computation (24BCE0556), parsing-table construction
with conflict detection (24BCI0103), and the predictive parser with
syntax-directed translation (24BDS0044) --- while the integration driver,
the test suite, the documentation and the review material were produced
jointly.

\section{LIMITATIONS}
The following limitations are stated as properties of the delivered
system, not as possibilities:
\begin{enumerate}
  \item \textbf{Fixed capacities.} All buffers are static arrays bounded
        by Table~\ref{tab:limits}; very large grammars or inputs are
        rejected safely rather than handled.
  \item \textbf{No grammar transformation.} Left-recursive or
        left-factoring-required grammars are reported NOT LL(1) with
        their conflicting cells; the tool does not rewrite them.
  \item \textbf{No lexer layer.} Input tokens are whitespace-separated
        words; multi-character terminals like \code{id} are just words,
        and no regular-expression specification is supported.
  \item \textbf{SDT shape rules.} The built-in translation covers the
        standard operator-encoding pattern; RHS shapes outside it
        translate as plain concatenation (a documented echo) rather than
        failing.
  \item \textbf{Trace and semantic caps.} Extremely long inputs truncate
        the printed trace (with a notice) or stop the translation with an
        error at the documented limits.
  \item \textbf{LL(1) only.} No conflict-resolution heuristics and no LR
        fallback; SLR/LALR behaviour is out of scope.
  \item \textbf{Implicit `$\$$' column.} The end-marker column is
        functional but not printed as part of the terminal row.
\end{enumerate}

\section{SCOPE FOR FUTURE WORK}
The architecture makes each enhancement a bounded, local change:
\begin{enumerate}
  \item \textbf{Grammar transformation stage:} automatic left-recursion
        elimination and left factoring before table construction, with a
        before/after diff shown to the user.
  \item \textbf{Richer diagnostics:} suggestions mode that names the
        offending alternative pair and the applicable LL(1) condition
        (the three-condition test) alongside the conflict report.
  \item \textbf{Parse-tree visualisation:} the parser's reduction history
        already encodes a leftmost derivation; rendering it as a tree
        would complete the classical toolset.
  \item \textbf{Additional translation schemes:} infix-to-prefix, syntax
        tree construction and typed TAC emission behind the existing
        event hook.
  \item \textbf{Formats and ergonomics:} multi-line alternatives,
        comments in grammar files, and a library mode exposing the
        modules for classroom experiments.
  \item \textbf{Dynamic storage normalization:} replacing the largest
        static buffers with arena allocation to lift the capacity limits
        while keeping the bounds-checking discipline.
\end{enumerate}



%% ======================= ANNEXURES =========================================

\appendix

\renewcommand{\chaptertitlename}{ANNEXURE}

\chapter{TEST CASES}
\label{annex:testcases}

This annexure lists the automated integration suite
(\code{testcases/run\_tests.sh}) in the standard test-case format. All
cases are executed by one suite run whose bottom line is printed after
the register; the individual execution date should be recorded in
Table~\ref{tab:setup} (Chapter 5) when you re-verify on your machine.

\section{TEST CASE TEMPLATE}
\begin{table}[h]
  \caption{Test case template (per-check)}
  \label{tab:tc_template}
  \small
  \begin{tabularx}{\textwidth}{p{3.2cm}X}
    \toprule
    \textbf{Field} & \textbf{Content} \\
    \midrule
    Test Case ID & Txx (Txxb: sub-run of the same grammar) \\
    Test Case Description & One-sentence intent of the scenario \\
    Test Case Type & Integration (black-box, via the CLI binary) \\
    Pre-conditions & Executable built; testcases/ tree present \\
    Input & Grammar file path and/or stdin strings \\
    Expected / Actual Result & Marker strings the suite greps for \\
    Pass/Fail & Status of the most recent executed run \\
    \bottomrule
  \end{tabularx}
\end{table}

\section{EXECUTED TEST CASE REGISTER}
\begin{longtable}{@{}p{0.9cm}p{2.5cm}p{3.2cm}>{\raggedright\arraybackslash}p{3.9cm}>{\raggedright\arraybackslash}p{1.5cm}@{}}
  \caption{Complete executed test-case register (26 checks)}
  \label{tab:tc_all}\\
  \toprule
  \textbf{ID} & \textbf{Description} & \textbf{Input} & \textbf{Expected
  marker} & \textbf{Status} \\
  \midrule
  \endfirsthead
  \toprule
  \textbf{ID} & \textbf{Description} & \textbf{Input} & \textbf{Expected
  marker} & \textbf{Status} \\
  \midrule
  \endhead
  \bottomrule
  \endfoot
  T01 & Demo grammar end-to-end & expression\_valid.txt & FIRST/FOLLOW
  exact sets; ``Grammar is LL(1)'' & PASS \\
  T02 & Epsilon propagation & epsilon\_propagation.txt & FIRST(S)$=$\{a
  ,b\}; nullable A,B & PASS \\
  T03 & Pure-epsilon grammar & pure\_epsilon.txt & FIRST(S)$=$\{\#\};
  LL(1) & PASS \\
  T03b & Epsilon input string & empty line on pure\_epsilon & INPUT
  ACCEPTED & PASS \\
  T04 & Conflict detection & non\_ll1\_conflict.txt & ``Conflict
  detected at M[S, a]''; NOT LL(1) & PASS \\
  T05 & Minimal grammar & simple.txt (S$\to$a) & FIRST(S)$=$\{a\};
  FOLLOW(S)$=$\{$\}$ & PASS \\
  T05b & Accept valid input & ``a'' & INPUT ACCEPTED & PASS \\
  T05c & Reject invalid input & ``b'' & SYNTAX ERROR; no entry [S,b] &
  PASS \\
  T06a & Undefined non-terminal & undefined\_nt.txt & rejected: ``never
  defined'' & PASS \\
  T06b & Missing arrow & missing\_arrow.txt & rejected: ``missing
  `-\!>' '' & PASS \\
  T06c & Double arrow & double\_arrow.txt & rejected: ``more than one
  arrow'' & PASS \\
  T06d & Invalid symbol & bad\_symbol.txt (a\$b) & rejected: ``invalid
  symbol'' & PASS \\
  T06e & Empty alternative & empty\_alt.txt & rejected: ``empty
  alternative'' & PASS \\
  T06f & End marker in grammar & end\_marker.txt & rejected: `$\$$'
  reserved & PASS \\
  T07 & Indirect recursion & indirect\_recursion.txt & correctly NOT
  LL(1) (M[A,c] clash) & PASS \\
  T08 & Valid input + trace & demo grammar, 3 strings & INPUT ACCEPTED
  with full trace & PASS \\
  T09 & Invalid input (mid) & ``id + * id'' & SYNTAX ERROR no entry [T,*] &
  PASS \\
  T09b & Invalid input (end) & ``id +'' & SYNTAX ERROR & PASS \\
  T10 & SDT operators & ``a + b * c'' & postfix ``a b c * +''; TAC t1,$t2$ &
  PASS \\
  T10b & SDT variant & ``id * id + id'' & postfix ``id id * id +''; TAC
  t1, t2 & PASS \\
  T11 & Compact arrow syntax & no\_space\_arrow.txt & accepted; LL(1) &
  PASS \\
  T12 & Nullable in RHS & optional\_symbol.txt & accepted; LL(1) & PASS \\
  T12b & Optional parsing & ``a b'' and ``b'' & both ACCEPTED & PASS \\
  T13 & Left recursion & left\_recursive.txt & NOT LL(1) reported & PASS \\
  \bottomrule
\end{longtable}
\noindent Final line of the executed run:
\begin{lstlisting}[style=outputcode]
  Tests passed: 26   failed: 0
\end{lstlisting}

\chapter{CODE}
\label{annex:code}

The complete source (2\,976 lines across 8 modules, plus headers, tests
and documentation) is delivered in the project repository. This annexure
reproduces the three excerpts that carry the intellectual core of the
system, with the files referenced. \emph{(The listing line numbers refer
to the delivered files.)}

\section{FIRST Fixed-Point Sweep (src/first.c --- excerpt)}
\begin{lstlisting}[language=C,label=lst:firstcode,caption={FIRST fixed-point iteration (abridged from src/first.c)}]
int computeAllFirstSets(Grammar *g, FirstCollection *out) {
    initializeFirstCollection(out, g);
    int changed = 1;
    while (changed) {                       /* least fixed point   */
        changed = 0;
        for (int p = 0; p < g->productionCount; ++p) {
            Production *pr = &g->productions[p];
            for (int a = 0; a < pr->rhsCount; ++a) {
                if (productionAltIsEpsilon(pr->rhs[a])) {
                    addEpsilonToFirstSet(out, p);   /* A -> #
                    changed |= 1;
                } else {
                    char sym[MAX_SYMBOL_LENGTH];
                    int nullable = 1;
                    int cnt = productionAltSymbolCount(pr->rhs[a]);
                    for (int i = 0; i < cnt; ++i) {
                        productionAltGetSymbol(pr->rhs[a], i, sym);
                        addSymbolsToFirstSet(out, p, sym, g);
                        if (!firstSetContainsEpsilon(...)) {
                            nullable = 0; break;    /* stop the chain */
                        }
                    }
                    if (nullable) addEpsilonToFirstSet(out, p);
                }
            }
        }
    }
    return 1;
}
\end{lstlisting}
(Identifiers are lightly abridged; the delivered file is normative.)

\section{FOLLOW Rules Sweep (src/follow.c --- excerpt)}
\begin{lstlisting}[language=C,label=lst:follcode,caption={FOLLOW fixed point: Rules 1--3 (abridged from src/follow.c)}]
int initializeFollowCollection(void);
/* Rule 1: $ belongs to FOLLOW(start). */
addTerminalToFollowSet(follows, startIndex, END_MARKER, g);

while (changed) {
    for (each production A -> alpha) {
        for (each symbol position i of alpha) {
            if (symbol is non-terminal B) {
                beta = suffix of alpha after B;
                firstOfSequence(beta, g, firsts, seq, &eps);
                /* Rule 2: FIRST(beta)\{#} into FOLLOW(B). */
                addSequenceSymbolsExceptEpsilon(...);
                /* Rule 3: beta nullable (or empty) => FOLLOW(A). */
                if (betaEmpty || eps)
                    mergeFollowInto(follows, B_index, A_index, g);
            }
        }
    }
}
\end{lstlisting}

\section{Predictive Parse Loop (src/parser.c --- excerpt)}
\begin{lstlisting}[language=C,label=lst:parsecodang,caption={Table-driven predictive loop (abridged from src/parser.c)}]
/* stack = "$ start"; tokens already end with "$". */
while (!accepted && !state.lastError) {
    tracePrintRow(&state, ...);
    if (++state.stepNumber > MAX_PARSE_STEPS)
        { lastError = "step limit exceeded"; break; }
    const char *X = topSymbol(&state);
    const char *a = curToken(&state);
    if (isEndMarker(X)) {
        if (isEndMarker(a)) accepted = 1;             /* accept   */
        else trailingGarbageError(&state, a);          /* $ vs a   */
    } else if (grammarIsTerminal(g, X)) {
        if (matchToken(X, a)) { fireShift(X); pop(&state); advance(&state); }
        else mismatchError(&state, X, a);
    } else {
        TableEntry e = tableLookup(table, X, a);
        if (isEmptyEntry(&e)) noTableEntryError(&state, X, a);  /* +expect */
        else { popNonTerminal(&state);
               pushRhsReversed(&state, e.lhsIndex, e.altIndex);
               recordReduceMarker(e.lhsIndex, e.altIndex); }
    }
}
\end{lstlisting}

\section{Grammar Corpora Sample}
The three representative corpora of the suite
(\code{testcases/valid/\allowbreak expression\_valid.txt},
\code{testcases/invalid/\allowbreak missing\_arrow.txt},
\code{testcases/edge/\allowbreak pure\_epsilon.txt}):
\begin{lstlisting}[style=cfg,label=lst:corpora,caption={Sample grammar corpora}]
# valid/expression_valid.txt
E  -> T Ep
Ep -> + T Ep | #
T  -> F Tp
Tp -> * F Tp | #
F  -> ( E ) | id

# invalid/missing_arrow.txt   (line 1 has no "->")
E T Ep

# edge/pure_epsilon.txt
S -> #
\end{lstlisting}


%% ======================= REFERENCES ========================================

\chapter{REFERENCES}
\label{page:references}
\begin{thebibliography}{9}
\setlength{\itemsep}{6pt}

\bibitem{lewis1968}
P.~M. Lewis, II and R.~E. Stearns, ``Syntax-directed transduction,''
\emph{Journal of the ACM}, vol.~15, no.~3, pp.~465--488, Jul.~1968,
doi:~10.1145/321466.321477.

\bibitem{rosenkrantz1970}
D.~J. Rosenkrantz and R.~E. Stearns, ``Properties of deterministic
top-down grammars,'' \emph{Information and Control}, vol.~17, no.~3,
pp.~226--256, 1970, doi:~10.1016/S0019-9958(70)90446-8.

\bibitem{aho2006compilers}
A.~V. Aho, M.~S. Lam, R.~Sethi, and J.~D. Ullman, \emph{Compilers:
Principles, Techniques, and Tools}, 2nd~ed. Boston, MA, USA:
Pearson/Addison-Wesley, 2006, ISBN~978-0-321-48681-3.

\bibitem{aho1972}
A.~V. Aho and J.~D. Ullman, \emph{The Theory of Parsing, Translation, and
Compiling}, vol.~1: \emph{Parsing}. Englewood Cliffs, NJ, USA:
Prentice-Hall, 1972.

\bibitem{gruene2008}
D.~Grune and C.~J.~H. Jacobs, \emph{Parsing Techniques: A Practical
Guide}, 2nd~ed. London, U.K.: Springer, 2008, ISBN~978-1-4419-1901-4.

\bibitem{johnson1975}
S.~C. Johnson, \emph{Yacc: Yet Another Compiler-Compiler}, Computing
Sciences Technical Report No.~32, Bell Laboratories, Murray Hill, NJ,
USA, Jul.~1975.

\bibitem{bison2021}
Free Software Foundation, \emph{GNU Bison: The Yacc-Compatible Parser
Generator}, ver.~3.8.1, Sep.~2021. [Online]. Available:
\url{https://www.gnu.org/software/bison/manual/}

\bibitem{parr2013}
T.~Parr, \emph{The Definitive ANTLR 4 Reference}. Raleigh, NC, USA:
Pragmatic Bookshelf, 2013, ISBN~978-1-934356-99-9.

\bibitem{isoC99}
ISO/IEC, \emph{Programming Languages --- C}, ISO/IEC 9899:1999, 2nd~ed.,
Geneva, Switzerland, Dec.~1999.

\bibitem{isoC11}
ISO/IEC, \emph{Information Technology --- Programming Languages --- C},
ISO/IEC 9899:2011, 3rd~ed., Geneva, Switzerland, Dec.~2011.

\end{thebibliography}

\end{document}
